% =============================================================================
% cap2_hardware/06_decisao_arquitetural.tex
% Seção 2.6 -- Arquitetura recomendada (v2.2 — terceira revisão: cenário único)
% =============================================================================
\section{Arquitetura recomendada}
\label{sec:decisao-arquitetural}

As cinco seções anteriores estabeleceram o problema material em camadas: dez pontos com geometria não padronizada, sete sem tomada elétrica e nove com insolação direta suficiente para alimentação fotovoltaica (Seção~\ref{sec:caracterizacao-pontos}); requisitos mínimos e desejáveis confrontados com a realidade operacional (Seção~\ref{sec:requisitos-tipologias}); conjunto de restrições de invólucro e geometria de captura (Seção~\ref{sec:acondicionamento-geometria}); regime \emph{offline-first} como princípio energético e comunicacional (Seção~\ref{sec:energia-conectividade}); e quatro tipologias arquiteturais hoje disponíveis, com seus respectivos \emph{trade-offs} (Seção~\ref{sec:tipologias-arquiteturais}). Esta seção articula essas camadas em decisão arquitetural concreta para o Campus 2 e fixa o quadro de parâmetros que o pipeline do Capítulo~\ref{cap:pipeline} recebe como entrada.

\subsection*{Síntese das restrições}

A leitura cruzada das seções anteriores mostra que o projeto opera em envelope estreito definido por seis vetores: heterogeneidade material entre pontos; escassez de tomadas em sete dos dez pontos, contrabalançada pela alta viabilidade solar em nove deles; instabilidade institucional da conectividade, que rebaixa a rede à condição de camada opcional e impõe \emph{offline-first} como princípio; biodiversidade não-alvo recorrente, que inviabiliza soluções baseadas apenas em detecção genérica de movimento e exige classificação multi-espécie; regime de baixa luminosidade noturna, que torna a iluminação por infravermelho próximo simultaneamente requisito funcional e item energético sensível; e contexto cooperativo da AEX Gatosdoc2, que recomenda baixa saliência visual do equipamento e integração às casinhas de reciclagem existentes. Esses seis vetores são as condições de contorno sob as quais qualquer aquisição é possível, e o pipeline do Capítulo~\ref{cap:pipeline} terá de aceitá-las como dadas.

\subsection*{Decisão arquitetural: tipologia única padronizada}

Em princípio, a heterogeneidade entre pontos abriria espaço para um \emph{arranjo arquitetural heterogêneo}, com tipologia escolhida ponto a ponto a partir do envelope material local. Essa opção, embora tecnicamente defensável, é incompatível com a escala operacional do projeto: um arranjo com quatro tipologias distintas exigiria quatro cadeias logísticas paralelas de aquisição, peça de reposição, treinamento, configuração e manutenção, em sistema mantido por equipe enxuta de voluntários sem dedicação técnica integral. A consequência prática é \textbf{a adoção de tipologia única padronizada para os dez pontos}, mesmo a custo de algum subaproveitamento de envelope em pontos com folga energética. A padronização viabiliza estoque rotativo de unidades equivalentes, substituição célere em caso de falha, procedimento único de configuração e treinamento mínimo da equipe de campo, alinhando a decisão arquitetural ao perfil de operação cooperativa documentado na Seção~\ref{sec:caracterizacao-pontos}.

Confrontados orçamento típico de projeto de graduação, disponibilidade no comércio nacional e opinião dos voluntários quanto à facilidade de manejo, a tipologia escolhida é a \textbf{câmera IP comercial doméstica adaptada, com gravação local em cartão microSD industrial, operação estritamente \emph{offline} e alimentação por \emph{powerbank}} (Seção~\ref{sec:tipologias-arquiteturais}, primeiro degrau). Quatro argumentos sustentam essa decisão. O primeiro é o \textbf{custo unitário}: câmeras IP de consumo doméstico são significativamente mais baratas que \emph{trail cameras} de média gama e que conjuntos SBC+NPU completos, o que viabiliza a cobertura simultânea dos dez pontos sob o orçamento disponível, com reserva para unidades de estoque rotativo. O segundo é a \textbf{disponibilidade no comércio nacional}: hardware, peças de reposição, fontes, \emph{powerbanks}, cartões microSD industriais e suportes mecânicos são facilmente encontrados, o que reduz o risco logístico e o tempo médio de recuperação de falhas em campo. O terceiro é a \textbf{simplicidade operacional para os voluntários}: o ciclo de campo se reduz à troca do \emph{powerbank} e à retirada periódica do cartão microSD, sem necessidade de painel fotovoltaico, controlador de carga ou rede sem fio dedicada por ponto --- procedimento já compatível com a rotina de alimentação semanal da colônia. O quarto é o \textbf{controle do nó de campo}: a substituição do \emph{firmware} de fábrica por uma imagem baseada no projeto OpenIPC~\cite{openipc2025}, quando o SoC suportar, permite operar em regime estritamente local, desabilitar rádios não utilizados, conter o consumo em modo ocioso e expor interface administrativa auditável, mitigando simultaneamente as fragilidades energética, de privacidade e de dependência de fornecedor associadas ao uso da câmera sem modificação.

\subsection*{Alternativas reservadas a testes e evoluções futuras}

A escolha por tipologia única não exclui as demais tipologias da escada tecnológica; reserva-as a finalidades específicas no horizonte de evolução do projeto. As \emph{trail cameras} alimentadas por baterias permanecem como referência comparativa para validação experimental nos pontos sem insolação adequada e em ciclos de campo nos quais se queira contrastar a taxa de \emph{ghost frames} e a qualidade noturna entre tipologias. Conjuntos SBC com NPU mantêm-se como referência aspiracional para etapas do pipeline em borda, podendo ser implantados em uma ou duas unidades-piloto para avaliar o ganho de descarte antecipado de quadros sem animal e de classificação multi-espécie embarcada, em preparação a futuras versões do sistema que exijam operação próxima ao tempo real. Os módulos baseados em microcontrolador (ESP32-CAM e similares) ficam reservados a redes densas de prova de conceito ou a pontos cuja perda do equipamento por vandalismo seja considerada aceitável. Em todos os casos, essas tipologias entram no projeto como \textbf{unidades experimentais paralelas}, não substituindo a tipologia padronizada nos dez pontos operacionais; sua função é fornecer evidência empírica para eventual migração arquitetural em versões futuras do sistema.

\subsection*{Consequências para o pipeline}

A adoção de câmeras IP comerciais em modo \emph{offline} como tipologia única tem três consequências metodológicas diretas para o Capítulo~\ref{cap:pipeline}. Em primeiro lugar, o pipeline pode pressupor, como caso de uso operacional, a chegada de quadros em formato comprimido (JPEG ou H.264) a partir de câmeras IP em modo \emph{offline}, sem metadados de pré-inferência embarcada, devendo concentrar capacidade de processamento integral em servidor. Em segundo lugar, a uniformidade da tipologia simplifica o estágio de normalização e calibração, que passa a tratar variações de geometria e iluminação entre pontos sem precisar absorver heterogeneidade adicional vinda do próprio sensor. Em terceiro lugar, a cadência de chegada de dados é determinada pelo ciclo logístico de retirada de cartões e troca de \emph{powerbanks} pelos voluntários, com periodicidade semanal ou quinzenal, o que estrutura o pipeline em regime de \textbf{lote} sobre janelas históricas e dispensa hipóteses sobre latência de transmissão. A presença sistemática de fauna não-alvo, herdada da Seção~\ref{sec:caracterizacao-pontos}, mantém a classificação multi-espécie como pré-requisito funcional da re-ID, agora aplicada a um regime visual de entrada uniforme entre pontos.

A discussão acima manteve-se no plano qualitativo. Os atributos de campo consolidados por ponto (energia, conectividade, incidência solar, fauna não-alvo, risco de vandalismo, entre outros) são reunidos na matriz P01--P10 do Apêndice~\ref{ap:matriz-pontos}, que serve como evidência de engenharia da decisão arquitetural aqui sintetizada. O detalhamento numérico do dimensionamento de \emph{powerbanks} e cartões microSD, das classes IP, dos ângulos de montagem e dos \emph{datasheets}, fora do foco deste trabalho centrado no software de visão computacional, é parte do projeto executivo subsequente.
